\section*{Case Studies: Learning Objectives Covered by Thesis}

For ESIPS substitution, this thesis is proposed to replace content from two elective
units: \textbf{AMME5520 Advanced Control and Optimisation} and \textbf{ELEC3305 Digital
Signal Processing}. The thesis provides complete coverage of all learning objectives from
each unit outline through the self-guided learning, implementation, and empirical
evaluation required to meet its research contributions.

\subsection*{Case Study 1: AMME5520 Advanced Control and Optimisation}

\subsubsection*{LO1. \emph{Approach research papers in a professional and research-
orientated manner, and conduct critical reviews of these papers}}
This learning objective is fully met through the thesis literature review (Chapter 3),
which critically evaluates over 30 papers across four domains: UAV DAA radar sensing,
GPU-accelerated radar processing, monopulse angle estimation under low SNR, and
multi-target tracking. For each work, assumptions are interrogated, experimental
limitations are identified (e.g.\ the absence of end-to-end pipeline validation in
Zaknoun et al.\ and the GPU benchmarking gap in Newmeyer et al.), and design decisions
are justified with reference to evidence. The review culminates in a structured
synthesis of three literature gaps that directly motivate the thesis contributions,
demonstrating the critical engagement with research central to this learning objective.

\subsubsection*{LO2. \emph{Implement simple path generation algorithms, controllers and
decision metrics for an autonomous system in order to meet specific mission objectives}}
This learning objective is met by implementing the estimation and decision components of
a multi-target tracking chain that is architecturally equivalent to the estimation layer
of any autonomous system. Specifically, the thesis implements: (i) a
\textbf{Kalman filter} for state prediction and update, including covariance propagation
through the predict--update cycle; (ii) \textbf{Mahalanobis distance gating}, which
applies an ellipsoidal decision boundary in state space to determine whether a detection
is plausibly associated with a given track; and (iii) a \textbf{Global Nearest Neighbour
(GNN)} data associator that resolves the assignment of detections to tracks as a minimum-
cost bipartite matching problem. Track initiation, promotion, and termination logic
provides mission-level decision-making over track lifetime. The tracker is validated on
simulated multi-target datasets across five noise regimes, achieving position errors
below 22\,m RMS across all conditions. This directly aligns with Weeks 07--08 (Kalman
filtering and extensions) and the unit's emphasis on state-based decision-making in
autonomous systems.

\subsubsection*{LO3. \emph{Understand a number of different path generation and control
algorithms implemented in autonomous systems and how they are linked to optimality
criteria, platform stability and vehicle constraints.}}
This learning objective is met through the comparative study of algorithmic alternatives
throughout the pipeline, framed in terms of optimality and constraint satisfaction.
The multi-target tracking stage investigates the GNN association algorithm under
varying clutter densities and SNR levels, examining how the gating threshold controls
the stability--sensitivity trade-off: too narrow a gate causes track fragmentation;
too wide a gate causes spurious associations. The analysis of CFAR variants
(CA-CFAR, GO-CFAR, SO-CFAR) under controlled noise conditions constitutes a
comparative study of detection algorithms linked to the false-alarm/sensitivity
optimality criterion, finding that CA-CFAR's larger training cell count produces
a more stable noise floor estimate and outperforms GO-CFAR and SO-CFAR in high-noise
conditions. Windowing evaluation across Hann, Hamming, Blackman, Kaiser, and Taylor
functions similarly frames each choice in terms of the resolution--sidelobe-attenuation
trade-off under SWaP constraints. This aligns with Weeks 02--03 (optimality-driven
algorithms) and Weeks 07--09 (estimation, uncertainty, robustness).

\subsubsection*{LO4. \emph{Understand how cost functions are used to define mission
objectives in a mathematical form, so that autonomous systems can make decisions about
their next action.}}
This learning objective is met by treating pipeline design as a formal optimisation
problem in which multiple competing objectives are balanced against hardware constraints.
The Mahalanobis distance used in gating is itself a cost function, quantifying
the statistical distance between a detection and a predicted track state in units of
standard deviations. The CFAR $\alpha$ coefficient is derived numerically by solving a
minimisation problem over the false alarm rate equation --- a bisection search that
finds the threshold value satisfying a formal probability target. An ablation study
systematically varies integration segment count and measures the resulting latency and
throughput, providing a cost--benefit surface over the design space. Parameter sweeps
over CFAR window size, guard cell count, and SNR are used to identify Pareto-optimal
operating points. This aligns with Week 01 (optimal control framing), Weeks 10--11
(optimisation/MPC), and the unit's emphasis on formal objective definition.

\subsubsection*{Coverage statement}
The thesis provides complete coverage of LO1--LO4 through: a critical multi-paper
literature review with structured gap synthesis (LO1); implementation and validation of
a Kalman filter, Mahalanobis gating, and GNN associator in a real-time multi-target
tracking system (LO2); comparative evaluation of detection and tracking algorithms
linked to formal optimality criteria and stability constraints (LO3); and pipeline
design as a formal multi-objective optimisation problem with numerical cost functions and
parameter sweep analysis (LO4).

\subsection*{Case Study 2: ELEC3305 Digital Signal Processing}

\subsubsection*{LO1. \emph{Describe signals using mathematical concepts such as
convolutions, transforms, spectral analyses, linear difference equations, filters,
correlation and covariance}}
This learning objective is fully met through the thesis's treatment of radar signal
processing. The dechirping process applies \textbf{matched filtering}, a convolution
between the received signal and a reference chirp, to compress pulse energy into a
narrow range bin. An \textbf{FFT} is then applied across fast time (per-chirp) to
convert the time-domain samples into a range-frequency map, and across slow time
(across chirps) to resolve Doppler. Five \textbf{window functions} --- Hann, Hamming,
Blackman, Kaiser, and Taylor --- are formulated and applied as element-wise
multiplications whose frequency-domain effect is analysed as a convolution with the
window's spectral response, quantifying the sidelobe attenuation and mainlobe width
trade-off for each. \textbf{Coherent integration} accumulates complex FFT outputs by
summation, exploiting phase coherence of genuine target returns while averaging out
incoherent noise. The Kalman filter's predict--update equations use \textbf{covariance
matrices} to propagate measurement and process uncertainty, and the Mahalanobis
distance is a covariance-normalised norm used for gating. These elements together
demonstrate mastery of convolution, transforms, spectral analysis, filtering, and
covariance as described in this learning objective.

\subsubsection*{LO2. \emph{Use software, such as MATLAB and Python, to develop code to
analyse and process signals including filtering, inverse filtering, spectral analyses,
resampling, time-frequency analyses}}
This learning objective is met through substantial software implementation across
multiple languages and platforms. The radar processing pipeline is implemented in
\textbf{C++ and CUDA}, including all DSP stages from windowing through FFT, coherent
integration, CFAR, clustering, monopulse angle sensing, and multi-target tracking.
A \textbf{Python} simulation framework generates synthetic radar datasets (DS1 and DS2)
and performs quantitative evaluation of detection and tracking metrics including recall,
position error RMS, and track fragmentation. A live visualisation harness renders FFT
waterfall plots, CFAR thresholds, cluster centroids, and 3D track histories in real
time. The pipeline achieves end-to-end latency below 50\,ms on an NVIDIA Xavier AGX,
with throughput of 9.8$\times$ the incoming data rate on a desktop and 3.8$\times$ on
the Xavier NX. CPU versus GPU implementations of CFAR are benchmarked across buffer
sizes and hardware configurations, directly addressing the software-based DSP
implementation emphasis of Weeks 02--12 practicals.

\subsubsection*{LO3. \emph{Identify linear signal models for time-invariant systems such
as convolutional models, sinusoidal and harmonic models, echo and reverberation models,
modulation models and time-frequency models.}}
This learning objective is met by modelling radar returns as \textbf{complex sinusoidal
signals} whose frequency encodes target range via the FMCW beat-frequency relationship,
and applying the \textbf{range--Doppler time-frequency model} to jointly represent range
and velocity. The synthetic dataset DS1 generates binary streams of complex sine waves
with additive Gaussian noise, providing a controlled sinusoidal signal model against
which FFT output can be analytically verified. The spectral leakage behaviour of each
window function is characterised through its frequency-domain response --- a
\textbf{convolutional model} linking the window's spectral shape to the observed
sidelobes in the range map. The effect of low SNR on signal model validity is studied
empirically: at $-10$\,dB SNR, individual monopulse angle detections become noisy, but
accumulating detections over time recovers unbiased estimates, consistent with the
linear model's noise averaging properties. This aligns with Week 02 (signal models),
Weeks 04--06 (DTFT/Z-transform/transfer functions), and Weeks 11--12 (sampling and
ADC/DAC considerations).

\subsubsection*{LO4. \emph{Design and evaluate signal processing systems.}}
This learning objective is directly and comprehensively met. The thesis designs an
end-to-end monopulse radar signal processing pipeline and evaluates each stage
systematically against detection and tracking performance metrics. Five window functions
are evaluated for sidelobe attenuation and its downstream effect on CFAR false alarm
rate, finding that the Taylor window calibrated to the CFAR noise floor at $-50$\,dB
achieves the highest recall in the noisiest conditions. Coherent versus non-coherent
integration are evaluated for SNR improvement, with coherent integration of 64 chirps
found to raise detection recall from near zero to 0.9 in the worst noise regime. Three
CFAR variants (CA-CFAR, GO-CFAR, SO-CFAR) are evaluated against false alarm rate and
detection sensitivity, and three CFAR implementations (naive CPU, prefix-sum CPU, GPU)
are benchmarked for latency and throughput across hardware platforms. An ablation study
identifies CFAR as the latency bottleneck of the tracking subsystem, contributing
1.5--5$\times$ more latency than the angle sensing stage. Monopulse angle sensing is
validated as an unbiased bearing and elevation estimator even at $-10$\,dB SNR, and the
full pipeline is verified on a live airborne field trial that geolocated static ground
targets at ranges of 410\,m, 603\,m, and 706\,m. This aligns with Weeks 04--10 and the
unit's system design emphasis.

\subsubsection*{Coverage statement}
The thesis provides complete coverage of LO1--LO4 through: mathematical treatment of
convolution, transforms, windowing, and covariance throughout the pipeline (LO1);
multi-language software implementation and benchmarking of the full DSP chain on real
embedded hardware (LO2); application and empirical validation of sinusoidal, echo, and
time-frequency signal models under controlled and real-world noise regimes (LO3); and
end-to-end design and quantitative evaluation of a real-time radar signal processing
system with live field trial validation (LO4).
